% Opciones originales: a. No requiere definir cuotas previas. | b. No usa selección aleatoria dentro de los grupos. | c. Es más preciso estadísticamente.
1 & El muestreo por cuotas se parece al estratificado, pero: & \textbf{b. No usa selección aleatoria dentro de los grupos.} & Ambos distinguen grupos, pero cuotas completa cantidades prefijadas mediante selección no aleatoria. El estratificado selecciona unidades al azar dentro de cada estrato; esa diferencia impide atribuir automáticamente a cuotas el error de un diseño probabilístico. \citep{statcanNoProb2021}.\\
\midrule
% Opciones originales: a. La población es homogénea. | b. No se dispone de información sobre la población. | c. Se desea asegurar la representación de subgrupos.
2 & El muestreo estratificado se aplica cuando: & \textbf{c. Se desea asegurar la representación de subgrupos.} & Se divide la población en estratos definidos y se extrae una muestra aleatoria en cada uno, lo que asegura presencia de los subgrupos previstos en el diseño. La homogeneidad dentro de estratos y diferencias entre ellos pueden mejorar la precisión. \citep{holmesBusiness2023,statcanProb2021}.\\
\midrule
% Opciones originales: a. El error de muestreo. | b. Una medida descriptiva de la muestra. | c. Una medida descriptiva de la población.
3 & ¿Qué representa el parámetro? & \textbf{c. Una medida descriptiva de la población.} & Un parámetro resume una característica numérica de la población completa, por ejemplo su media o proporción. La medida calculada con los datos de una muestra es un estadístico, que puede utilizarse para estimar el parámetro. \citep{wackerly2010}.\\
\midrule
% Opciones originales: a. La población está dispersa geográficamente. | b. La población es muy homogénea. | c. Se requiere seleccionar individuos específicos por criterio.
4 & El muestreo por conglomerados se usa cuando: & \textbf{a. La población está dispersa geográficamente.} & La selección de grupos naturales, como áreas o sucursales, permite concentrar la recolección y reducir desplazamientos cuando las unidades están geográficamente dispersas. El diseño selecciona conglomerados al azar; elegir individuos por criterio corresponde a un procedimiento no probabilístico. \citep{statcanProb2021}.\\
\midrule
% Opciones originales: a. Cada k-ésimo elemento después de un punto inicial aleatorio. | b. Elementos elegidos por juicio del investigador. | c. Solo individuos con características comunes.
5 & En el muestreo sistemático se selecciona: & \textbf{a. Cada k-ésimo elemento después de un punto inicial aleatorio.} & El procedimiento fija un intervalo de selección k, elige aleatoriamente el arranque y toma cada k-ésima unidad del listado. Debe revisarse que la periodicidad del marco no coincida con el intervalo para evitar distorsiones. \citep{holmesBusiness2023,statcanProb2021}.\\
\midrule
% Opciones originales: a. Elaborar la hipótesis nula | b. Seleccionar el marco muestral. | c. Definir objetivos.
6 & ¿Cuál de los siguientes NO es un paso en una encuesta de muestreo? & \textbf{a. Elaborar la hipótesis nula} & Definir objetivos y determinar el marco son componentes del diseño de una encuesta. Formular una hipótesis nula corresponde específicamente a una prueba de hipótesis y no constituye un paso indispensable de toda encuesta, que también puede tener fines descriptivos o de estimación. \citep{holmesBusiness2023}.\\
\midrule
% Opciones originales: a. El margen de error. | b. La proporción esperada. | c. El valor crítico de la distribución normal para el nivel de confianza
7 & En la fórmula n = \((Z\cdot\sigma/E)^2\), el valor Z representa: & \textbf{c. El valor crítico de la distribución normal para el nivel de confianza} & Para estimar una media con desviación poblacional conocida, Z es el cuantil normal asociado al nivel de confianza elegido; por ejemplo, 1.96 para un intervalo bilateral del 95 \%. E representa el margen de error y σ la desviación estándar poblacional. El tamaño se redondea hacia arriba. \citep{holmesBusiness2023}.\\
\midrule
% Opciones originales: a. Aplicar una encuesta a toda la población nacional. | b. Auditar una muestra de transacciones financieras. | c. Evaluar todos los empleados de una empresa.
8 & Un ejemplo de aplicación del muestreo en administración es: & \textbf{b. Auditar una muestra de transacciones financieras.} & La auditoría examina una parte de las transacciones para obtener evidencia sobre el conjunto delimitado, con un diseño y riesgo de muestreo explícitos. Observar a todos los integrantes de una población constituye un censo, por lo que las otras opciones no describen muestreo. \citep{holmesBusiness2023}.\\
\midrule
% Opciones originales: a. La selección de grupos enteros. | b. La selección por conveniencia. | c. Que todos los individuos tienen la misma probabilidad de ser elegidos.
9 & El muestreo aleatorio simple se caracteriza por: & \textbf{c. Que todos los individuos tienen la misma probabilidad de ser elegidos.} & En una muestra aleatoria simple todas las muestras posibles del mismo tamaño tienen igual probabilidad; de ello resulta la misma probabilidad de inclusión para cada unidad. Esa igualdad individual es una característica, aunque por sí sola no define completamente el diseño. \citep{holmesBusiness2023}.\\
\midrule
% Opciones originales: a. 0.5 | b. 1.0 | c. 0.25
10 & Cuando no se conoce la proporción poblacional (p), se usa el valor: & \textbf{a. 0.5} & En la planificación del tamaño muestral para una proporción sin estimación previa se utiliza p = 0.5 porque maximiza p(1-p) = 0.25. Así se obtiene el tamaño conservador bajo los supuestos de la fórmula; no significa que la proporción real sea necesariamente 50 \%. \citep{holmesBusiness2023}.\\
\midrule
% Opciones originales: a. Los datos se pierden durante el análisis. | b. La muestra no representa bien a la población por el método de selección. | c. Las preguntas influyen en la respuesta.
11 & El sesgo de selección ocurre cuando: & \textbf{b. La muestra no representa bien a la población por el método de selección.} & La selección favorece o excluye sistemáticamente unidades con características relacionadas con el resultado, de modo que la muestra representa inadecuadamente a la población objetivo. La redacción que induce respuestas corresponde a sesgo de medición; la pérdida de datos durante el análisis es otro problema. \citep{holmesBusiness2023,bradley2021}.\\
\midrule
% Opciones originales: a. Una medida descriptiva obtenida de la muestra. | b. Una medida descriptiva de la población. | c. Un error sistemático del muestreo.
12 & ¿Qué es el estadístico? & \textbf{a. Una medida descriptiva obtenida de la muestra.} & El estadístico es una función de los datos muestrales, como la media muestral o la proporción observada. Describe la muestra y puede actuar como estimador de un parámetro de la población; no es por definición un error. \citep{wackerly2010}.\\
\midrule
% Opciones originales: a. El sesgo del investigador. | b. Los errores de cálculo en la encuesta. | c. La diferencia entre el valor de la muestra y el parámetro poblacional.
13 & El error de muestreo se refiere a: & \textbf{c. La diferencia entre el valor de la muestra y el parámetro poblacional.} & Es la diferencia entre un estadístico muestral y el parámetro que estima atribuible a observar una muestra en lugar de toda la población. La variabilidad entre muestras no equivale al sesgo sistemático ni a errores de cálculo o registro, que son errores ajenos al muestreo. \citep{wackerly2010,holmesBusiness2023}.\\
\midrule
% Opciones originales: a. Un parámetro estadístico. | b. Todo el conjunto de individuos disponibles. | c. Un subconjunto representativo de la población
14 & La muestra es: & \textbf{c. Un subconjunto representativo de la población} & Es la opción adecuada porque una muestra es un subconjunto de unidades de la población que se observa para estudiarla. La representatividad es un objetivo que debe evaluarse mediante cobertura, selección y respuesta: pertenecer a una muestra no garantiza por sí mismo representar bien a la población. \citep{holmesBusiness2023,bradley2021}.\\
\midrule
% Opciones originales: a. Un método para estudiar toda la población. | b. Un procedimiento para eliminar datos atípicos. | c. Un proceso de selección de un subconjunto representativo de la población.
15 & ¿Qué es el muestreo? & \textbf{c. Un proceso de selección de un subconjunto representativo de la población.} & El muestreo establece cómo seleccionar una parte de la población para obtener información sobre ella con menor costo que un censo. Se busca una representación adecuada, pero esta depende del marco y del diseño y no está garantizada automáticamente. Eliminar valores atípicos es una operación distinta. \citep{holmesBusiness2023}.\\
\midrule
% Opciones originales: a. No disminuye con el tamaño de la muestra y es sistemático. | b. Disminuye al aumentar el tamaño de la muestra. | c. Es un error aleatorio sin consecuencias.
16 & El sesgo se diferencia del error de muestreo porque: & \textbf{a. No disminuye con el tamaño de la muestra y es sistemático.} & El sesgo es una desviación sistemática causada, por ejemplo, por cobertura o selección inadecuadas. Aumentar n reduce normalmente la variabilidad aleatoria bajo un diseño adecuado, pero no elimina por sí solo el mecanismo que origina el sesgo; este puede persistir en muestras enormes. \citep{bradley2021,holmesBusiness2023}.\\
\midrule
% Opciones originales: a. La división proporcional de la población. | b. Un algoritmo aleatorio. | c. La experiencia o criterio del investigador.
17 & El muestreo por juicio se basa en: & \textbf{c. La experiencia o criterio del investigador.} & La selección depende del conocimiento y propósito del investigador. Es no probabilística: el criterio experto no establece por sí solo probabilidades de inclusión conocidas ni garantiza representatividad. \citep{statcanNoProb2021}.\\
\midrule
% Opciones originales: a. Poblaciones difíciles de localizar o pequeñas. | b. Poblaciones grandes y homogéneas. | c. Muestreos de control de calidad.
18 & El muestreo bola de nieve se utiliza especialmente en: & \textbf{a. Poblaciones difíciles de localizar o pequeñas.} & Los contactos iniciales remiten a otros participantes y facilitan acceder a grupos poco visibles. La dificultad de localización es decisiva; el tamaño pequeño por sí solo no exige este método. \citep{statcanNoProb2021}.\\
\midrule
% Opciones originales: a. Solo los individuos que responden una encuesta. | b. El conjunto completo de elementos con una característica común | c. Un grupo pequeño seleccionado al azar.
19 & ¿Qué se entiende por población en estadística? & \textbf{b. El conjunto completo de elementos con una característica común} & La población comprende la totalidad de unidades que cumplen los criterios definidos para el estudio, con delimitación de ámbito, tiempo y características. Los respondientes son solo las unidades de las que se obtuvo información y no necesariamente toda la población objetivo. \citep{holmesBusiness2023}.\\
\midrule
% Opciones originales: a. Mediante estratificación previa. | b. Aleatoriamente de un marco muestral. | c. Por accesibilidad o disponibilidad.
20 & En el muestreo por conveniencia los participantes se eligen: & \textbf{c. Por accesibilidad o disponibilidad.} & Se selecciona a quienes resultan accesibles o disponibles, sin sorteo desde un marco. Esa facilidad puede excluir segmentos y limita la generalización de los resultados. \citep{holmesBusiness2023,statcanNoProb2021}.\\
